(1) 在热源瞬时温度 $T$ 时可逆机效率为 $\eta=1-T_0/T$，热源放热量为 $-C\,dT>0$。因此微元功 $dW=(1-T_0/T)(-C\,dT)$，并有环境吸热 $dQ_0=-CT_0\,dT/T$。这一计算默认每一步可逆、热机两端温差趋于无穷小；若有有限温差换热则实际功严格更少。
(2) 对 $T$ 自 $T_i$ 降至 $T_0$ 积分，$W_{\max}=C\int_{T_0}^{T_i}(1-T_0/T)\,dT=C[T_i-T_0-T_0\ln(T_i/T_0)]$。令 $x=T_i/T_0>1$，括号为 $T_0(x-1-\ln x)>0$，符合热机输出非负。与单纯可供热量 $C(T_i-T_0)$ 相比，必有一部分热传给环境以满足可逆熵平衡。
(3) 若直接与环境接触，热源内能减少 $C(T_i-T_0)$ 全部进入环境，熵增分别为 $\Delta S_{\rm hot}=C\ln(T_0/T_i)$、$\Delta S_{\rm env}=C(T_i-T_0)/T_0$。故 $\Delta S_{\rm gen}=C[(T_i-T_0)/T_0-\ln(T_i/T_0)]>0$。相乘得 $T_0\Delta S_{\rm gen}=W_{\max}$，表示直接散热毁掉的火用。
(4) 设 $\epsilon=\delta T/T_0\ll1$，用 $\ln(1+\epsilon)=\epsilon-\epsilon^2/2+\epsilon^3/3-\cdots$，得到 $W_{\max}=C\delta T^2/(2T_0)-C\delta T^3/(3T_0^2)+O(\delta T^4)$。首项为温差的二阶而非一阶，是因为近乎平衡时卡诺效率本身只有 $O(\delta T/T_0)$。热量的全部量不能转换为功。
